\begin{table}[ht]
\centering
\caption{Router output on all 30 challenge prompts (descriptive only)}
\label{tab:challenge_router_detail}
\begin{tabular}{lllllll}
\toprule
prompt\_id & challenge\_type & predicted\_capability\_id & fallback\_detected & api\_call\_succeeded & returned\_models & reason \\
\midrule
chal\_out\_of\_scope\_mixed\_01 & out\_of\_scope\_mixed & hello\_world & False & True & ['gpt-5.4-2026-03-05'] & The user asks for weather information, which does not match any analytical capability provided. Among the allowed options, only the generic demo capability is the closest fallback. \\
chal\_out\_of\_scope\_mixed\_02 & out\_of\_scope\_mixed & hello\_world & False & True & ['gpt-5.4-2026-03-05'] & The user requests creative writing, which does not match the analytical dataset capabilities. Among the allowed options, hello\_world is the only non-analytic fallback/demo capability. \\
chal\_out\_of\_scope\_mixed\_03 & out\_of\_scope\_mixed & hello\_world & False & True & ['gpt-5.4-2026-03-05'] & The user requests a simple language translation, which does not match any analysis/modeling capability. Among the allowed options, the only non-analytic fallback/demo capability is hello\_world. \\
chal\_out\_of\_scope\_mixed\_04 & out\_of\_scope\_mixed & linear\_regression & False & True & ['gpt-5.4-2026-03-05'] & The actionable analytical part of the request is to fit a regression of outcome on treatment in study.csv, which matches linear\_regression. The email request is outside the allowed capabilities, so select the single best matching capability for the regression task. \\
chal\_out\_of\_scope\_mixed\_05 & out\_of\_scope\_mixed & hello\_world & False & True & ['gpt-5.4-2026-03-05'] & The user asks for a general cooking recipe, which does not match any analytical/statistical capability. Among the allowed options, hello\_world is the only non-analytic fallback/demo capability. \\
chal\_out\_of\_scope\_mixed\_06 & out\_of\_scope\_mixed & summary\_stats & False & True & ['gpt-5.4-2026-03-05'] & The request includes summarizing outcomes.csv, which matches dataset summary statistics/EDA. Flight booking is unsupported, so choose the single applicable capability for the actionable data task. \\
chal\_out\_of\_scope\_mixed\_07 & out\_of\_scope\_mixed & hello\_world & False & True & ['gpt-5.4-2026-03-05'] & The user requests code review for a Python function, which does not match any analysis/modeling capability provided. Among the allowed options, hello\_world is the only non-analytic fallback. \\
chal\_out\_of\_scope\_mixed\_08 & out\_of\_scope\_mixed & hello\_world & False & True & ['gpt-5.4-2026-03-05'] & The user asks a general knowledge question that does not match any analytical capability; among the allowed options, hello\_world is the only non-analytic fallback. \\
chal\_out\_of\_scope\_mixed\_09 & out\_of\_scope\_mixed & causal\_ate & False & True & ['gpt-5.4-2026-03-05'] & The request contains two tasks, but only one capability can be chosen. Among the allowed capabilities, only estimating the average treatment effect of a drug on an outcome in study.csv matches a listed capability. No reminder-setting capability is available. \\
chal\_out\_of\_scope\_mixed\_10 & out\_of\_scope\_mixed & hello\_world & False & True & ['gpt-5.4-2026-03-05'] & The user requests investment advice, which is outside the scope of the available data-analysis and causal-inference capabilities. Since no listed capability fits financial advice, the only non-domain-specific option is hello\_world. \\
chal\_underspecified\_01 & underspecified & causal\_ate & False & True & ['gpt-5.4-2026-03-05'] & The user asks to analyze the effect of a treatment on outcomes, which is a general causal treatment-effect request best matched to estimating an average treatment effect rather than regression, EDA, or survival-curve-specific analysis. \\
chal\_underspecified\_02 & underspecified & causal\_ate & False & True & ['gpt-5.4-2026-03-05'] & The user asks whether the treatment worked, which implies estimating an overall treatment-versus-control causal effect. This is a standard average treatment effect question, and there is no explicit indication that the outcome is binary, continuous-model fitting, survival curves, or simple EDA. \\
chal\_underspecified\_03 & underspecified & causal\_ate & False & True & ['gpt-5.4-2026-03-05'] & The user asks for the impact of a drug in a clinical trial dataset, which is a general treatment-versus-control causal effect question best matched to estimating an average treatment effect. \\
chal\_underspecified\_04 & underspecified & summary\_stats & False & True & ['gpt-5.4-2026-03-05'] & The request is ambiguous and only asks to compare two treatment groups over the study period without explicitly requesting a causal effect estimate, regression model, or survival curves. The safest fit is exploratory dataset comparison/EDA on the CSV. \\
chal\_underspecified\_05 & underspecified & causal\_ate & False & True & ['gpt-5.4-2026-03-05'] & The user requests a generic causal analysis without specifying binary-outcome specialized estimation or survival curves, so the general-purpose average treatment effect method is the best fit. \\
chal\_underspecified\_06 & underspecified & causal\_ate & False & True & ['gpt-5.4-2026-03-05'] & The user asks whether drug X makes a difference for patients versus not taking it, which is a standard treatment-versus-control causal effect question seeking an overall effect estimate. There is no indication of survival curves, binary-outcome specialized estimation, regression-only modeling, or simple EDA. \\
chal\_underspecified\_07 & underspecified & causal\_ate & False & True & ['gpt-5.4-2026-03-05'] & The user asks to adjust for confounding and interpret what the treatment group data shows, which most closely matches a general causal treatment-versus-control effect estimation request. There is no explicit mention of binary outcome, regression modeling, survival curves, or simple EDA. \\
chal\_underspecified\_08 & underspecified & causal\_ate & False & True & ['gpt-5.4-2026-03-05'] & The user asks for the difference between treated and untreated groups, which implies estimating an overall treatment-versus-control effect. This is a standard causal effect request, not explicitly survival curves, binary-outcome specialized EDRIP, regression modeling, or simple summary statistics. \\
chal\_underspecified\_09 & underspecified & causal\_ate & False & True & ['gpt-5.4-2026-03-05'] & The user requests a single adjusted causal effect estimate from treatment\_data.csv. This matches general-purpose average treatment effect estimation with covariate adjustment, not survival curves, regression-only prediction, or the specialized binary-outcome EDRIP method since no binary-outcome emphasis is given. \\
chal\_underspecified\_10 & underspecified & causal\_ate & False & True & ['gpt-5.4-2026-03-05'] & The user asks whether an intervention helped and wants baseline differences between groups accounted for, which is a standard causal treatment-versus-control effect estimation request. They do not specify a binary outcome, regression modeling task, or survival curves, so the general-purpose adjusted average treatment effect method is the best fit. \\
chal\_unsupported\_method\_01 & unsupported\_method & linear\_regression & False & True & ['gpt-5.4-2026-03-05'] & The request is for estimating the effect of statin uptake on a continuous outcome (cholesterol). There is no instrumental-variables-specific capability available, so the closest available modeling option for a continuous outcome is linear\_regression. \\
chal\_unsupported\_method\_02 & unsupported\_method & causal\_ate & False & True & ['gpt-5.4-2026-03-05'] & The request is for estimating a causal effect of eligibility at a cutoff on an outcome. There is no dedicated regression discontinuity capability available, and the closest allowed option for treatment effect estimation is the general-purpose causal\_ate tool. \\
chal\_unsupported\_method\_03 & unsupported\_method & logistic\_regression & False & True & ['gpt-5.4-2026-03-05'] & The request is to estimate a difference-in-differences model on readmission rates, which are typically binary. Among the available options, logistic\_regression is the best fit for modeling a binary outcome using treated/control and pre/post indicators (and their interaction) as covariates. No dedicated DID capability is available. \\
chal\_unsupported\_method\_04 & unsupported\_method & causal\_ate & False & True & ['gpt-5.4-2026-03-05'] & The request is to estimate the effect of a policy intervention on an outcome using panel data. Synthetic control is a causal effect estimation task, and among the available options, causal\_ate is the closest general-purpose capability for estimating an overall treatment effect. Other options are for regression, EDA, binary-specific methods, or survival curves. \\
chal\_unsupported\_method\_05 & unsupported\_method & causal\_ate & False & True & ['gpt-5.4-2026-03-05'] & The request is for a causal effect decomposition, but no mediation-specific capability is available. Among the allowed options, causal\_ate is the closest causal inference tool for estimating treatment effects; the other options are for regression, EDA, binary specialized estimation, survival curves, or a demo. \\
chal\_unsupported\_method\_06 & unsupported\_method & causal\_ate & False & True & ['gpt-5.4-2026-03-05'] & The request is for causal treatment effect estimation. Although it specifically asks for a causal forest and heterogeneous treatment effects, no dedicated HTE/causal forest capability is available. Among the allowed options, causal\_ate is the closest causal inference capability for estimating treatment effects; regression, summary, binary-specific, and survival-curve tools are less appropriate. \\
chal\_unsupported\_method\_07 & unsupported\_method & causal\_ate & False & True & ['gpt-5.4-2026-03-05'] & The request is for causal effect estimation of treatment timing on mortality, not summary statistics, regression-only modeling, binary-outcome specialized EDRIP, or survival curve visualization. Among the available options, causal\_ate is the closest general-purpose causal inference capability, even though the requested parametric g-formula for time-varying treatment is more specialized than the listed tools. \\
chal\_unsupported\_method\_08 & unsupported\_method & causal\_ate & False & True & ['gpt-5.4-2026-03-05'] & The user requests estimation of the causal effect of an exposure on an outcome. Among the available options, causal\_ate is the closest general-purpose causal inference capability for estimating an overall treatment effect. Other options are for binary-outcome specialized estimation, regression, summary statistics, or survival curves, and there is no dedicated front-door capability available. \\
chal\_unsupported\_method\_09 & unsupported\_method & causal\_ate & False & True & ['gpt-5.4-2026-03-05'] & The user requests estimation of a matched-sample average treatment effect using propensity score matching. Among the available capabilities, causal\_ate is the general-purpose causal effect estimation option for a single overall treatment effect; the other options are for regression, EDA, binary specialized EDRIP, survival curves, or demo only. \\
chal\_unsupported\_method\_10 & unsupported\_method & causal\_ate & False & True & ['gpt-5.4-2026-03-05'] & The request is for causal effect estimation of a treatment/exposure using inverse probability weighting in a marginal structural model. Among the available options, causal\_ate is the closest general-purpose causal inference capability for estimating an overall treatment effect; the survival-curves tool is only for explicit survival curve requests, and the regression/EDA tools do not match the causal MSM request. \\
\bottomrule
\end{tabular}
\end{table}
